\begin{remark}\label{rl-0007}\label{rmk:description-of-homotopy-fixed-stack}
  Here is a slightly more explicit description of the points of $\mathcal X^{G}$ over a scheme $U \in \Sch/S$. An object of $\mathcal X^{G}(U)$ consists of a $G$-equivariant morphism $(x,\alpha):\iota(U)\to X$. The equivariance diagram in this case is
  \begin{center}
    \begin{tikzcd}[column sep={1.5cm,between origins},
      row sep={1.5cm,between origins}]
      G\times U\arrow[r, "\pr_2"] \arrow[d, swap, "\id_G \times x"] & U \arrow[d, "x"] \\
      G\times \mathcal X\arrow[r, "\mu"] \arrow[ur, Rightarrow, "\alpha"]& \mathcal X
    \end{tikzcd}
  \end{center}
  since $\iota(U)$ is simply $U$ with the trivial $G$-action. The morphism $x$ is an object of \(\mathcal X\) by the Yoneda lemma and the 2-commutativity implies that for all $g \in G(U)$ we get an isomorphism $\alpha_g:g\cdot x\xrightarrow{\sim}x$. Similarly, given a pair $(x,\{\alpha_g\}_{g\in G(U)})$ consisting of an object $x \in \mathcal X(U)$ and a collection of isomorphisms $\alpha_g:gx\xrightarrow{\sim}x$, we obtain a $G$-equivariant map $\iota(U)\to \mathcal X$. Thus the objects of $\mathcal X^{G}(U)$ are precisely pairs $(x,\{\alpha_g\}_{g\in G(U)})$ where $x \in \mathcal X(U)$ and the $\alpha_g:g\cdot x\xrightarrow{\sim} x$ are isomorphisms satisfying the cocycle condition \(\alpha_{gh} = \alpha_g\circ g\cdot \alpha_h\) pictured by
  \begin{equation}\label{eqn:cocycle-condition-for-fixed-stack}
    \begin{tikzcd}[column sep={1.5cm,between origins},
      row sep={1.5cm,between origins}]
      & gx \arrow[dr,"\alpha_g"] & \\
      (gh)x\arrow[ur, "g\cdot \alpha_h"] \arrow[rr, "\alpha_{gh}"] & & x
    \end{tikzcd}
  \end{equation}
  for all $g, h\in G(U)$. The morphism \(\eta:\mathcal X^G\to \mathcal X\) simply forgets the data of \(\alpha\).

  In particular, when \(G\) acts trivially on \(\mathcal X\), \(gx = x\) for all \(x\in \mathcal X(U)\) and the cocycle condition reduces to the requirement that \(\alpha_{gh} = \alpha_g\circ \alpha_h\); i.e. that \(\alpha:G\to \underline{\Aut}(x)\) is a group homomorphism. Thus, for the trivial action, the fiber of \(x\in \mathcal X(U)\) over \(\eta\) is the sheaf \(\underline{\Hom}_{S-gp}(G, \Aut(x))\).
\end{remark}
